% =================================================
% Ученический рабочий лист — скорректированная версия
% =================================================
\worksheettitle
  {Логарифмическое неравенство}
  {Степени логарифма, модуль и две последовательные замены}
\studentnameblock

\begin{checkpointbox}[Входная диагностика: 4 минуты]
  Выполните задания без подробного решения.
  \begin{enumerate}[label=\textbf{\arabic*.},itemsep=0.35em]
    \item Вычислите:
      \[
        \lg^2 100=\answerblank[24mm].
      \]
    \item Найдите область допустимых значений:
      \[
        \lg((x-1)^2),\qquad x\in\answerblank[52mm].
      \]
    \item Решите неравенство:
      \[
        y^2\ge9,\qquad y\in\answerblank[60mm].
      \]
    \item Решите неравенство:
      \[
        x^2\ge16,\qquad x\in\answerblank[60mm].
      \]
  \end{enumerate}
\end{checkpointbox}

\begin{notebox}[Как использовать диагностику]
  Это не оценочная работа. Отметьте задания, которые вызвали затруднение,
  чтобы понять, какой переход нужно повторить перед основной задачей.
\end{notebox}

\begin{problembox}
  Решите неравенство
  \[
    \lg^4((x^2-4)^2)-\lg^2((x^2-4)^4)\ge192.
  \]
\end{problembox}

\begin{notebox}[Сначала договоримся о записи]
  Верхний индекс у обозначения функции относится к \textbf{значению функции}:
  \[
    \lg^2 A=(\lg A)^2,
    \qquad
    \lg^4 A=(\lg A)^4.
  \]
  Поэтому исходное неравенство означает
  \[
    \bigl[\lg((x^2-4)^2)\bigr]^4
    -\bigl[\lg((x^2-4)^4)\bigr]^2\ge192.
  \]
\end{notebox}

\section{Учимся читать похожие записи}

\begin{fillbox}[Что означает число \(2\) в каждой записи?]
  Заполните таблицу: прочитайте каждую запись и укажите роль числа \(2\).
  \begin{center}
    \renewcommand{\arraystretch}{1.65}
    \begin{tabularx}{0.96\textwidth}{@{}
      >{\centering\arraybackslash}p{0.18\textwidth}
      >{\raggedright\arraybackslash}p{0.39\textwidth}
      >{\raggedright\arraybackslash}X@{}}
      \toprule
      \rowcolor{TableHeader}
      Запись & Как прочитать & Что означает число \(2\) \\
      \midrule
      \(\lg^2 A\)
        & \answerblank[48mm]
        & \answerblank[43mm] \\
      \addlinespace[0.35em]
      \(\lg(A^2)\)
        & \answerblank[48mm]
        & \answerblank[43mm] \\
      \addlinespace[0.35em]
      \(\log_2 A\)
        & \answerblank[48mm]
        & \answerblank[43mm] \\
      \bottomrule
    \end{tabularx}
  \end{center}
\end{fillbox}

\begin{notebox}[Самопроверка]
  Сравните свои ответы со схемой. Если обнаружили неточность, исправьте её
  другим цветом.
  \notationtrianglefigure
\end{notebox}


\section{Область допустимых значений}

Оба аргумента логарифмов должны быть строго положительными. Заполните цепочку:
\[
  (x^2-4)^2>0
  \quad\Longleftrightarrow\quad
  x^2-4\ne\answerblank[12mm]
  \quad\Longleftrightarrow\quad
  x\ne\answerblank[16mm],\;x\ne\answerblank[16mm].
\]

Нужно ли отдельно решать условие \((x^2-4)^4>0\)?
\choicebox\ да \qquad \choicebox\ нет.

Объясните ответ:
\writinglines{2}

\begin{checkpointbox}[ОДЗ]
  \[
    x\in\answerblank[80mm].
  \]
\end{checkpointbox}

\begin{reflectionbox}[Типичная ошибка]
  Почему рассуждение «аргумент является квадратом, значит он всегда положителен»
  неверно?
  \writinglines{2}
\end{reflectionbox}

\section{Почему появляется модуль}

\stepheading{Шаг 1}{Замечаем проблему на конкретном значении}

Выражение \(x^2-4\) может быть отрицательным. Подставьте \(x=0\):
\[
  x^2-4=\answerblank[16mm],
  \qquad
  (x^2-4)^2=\answerblank[16mm]>0.
\]
Какое выражение определено при \(x=0\)?
\[
  \square\; \lg((x^2-4)^2)
  \qquad
  \square\; 2\lg(x^2-4).
\]
Объясните, почему второе выражение может не иметь смысла:
\writinglines{1}

\stepheading{Шаг 2}{Сравниваем положительное и отрицательное значения}

\begin{center}
  \renewcommand{\arraystretch}{1.55}
  \begin{tabularx}{0.94\textwidth}{@{}
    >{\centering\arraybackslash}p{0.12\textwidth}
    >{\centering\arraybackslash}p{0.22\textwidth}
    >{\centering\arraybackslash}p{0.33\textwidth}
    >{\centering\arraybackslash}X@{}}
    \toprule
    \rowcolor{TableHeader}
    \(z\) & \(\lg(z^2)\) & Определено ли \(2\lg z\)? & \(2\lg|z|\) \\
    \midrule
    \(10\)  & \(2\) & да  & \(2\) \\
    \addlinespace[0.25em]
    \(-10\) & \answerblank[13mm] & \answerblank[18mm] & \answerblank[13mm] \\
    \bottomrule
  \end{tabularx}
\end{center}

Какая формула работает и при положительном, и при отрицательном \(z\)?
\[
  \square\; \lg(z^2)=2\lg z
  \qquad
  \square\; \lg(z^2)=2\lg|z|.
\]

\stepheading{Шаг 3}{Выводим общее правило}

Если \(z\ne0\), то \(|z|>0\). Дополните вывод:
\[
  z^2=\bigl(\answerblank[16mm]\bigr)^2,
\]
\[
  \lg(z^2)
  =\lg\bigl((\answerblank[16mm])^2\bigr)
  =2\lg\answerblank[16mm],
\]
\[
  \lg(z^4)
  =\lg\bigl((\answerblank[16mm])^4\bigr)
  =4\lg\answerblank[16mm].
\]

\begin{notebox}[Почему модуль необходим]
  Модуль заменяет число \(z\) положительным числом \(|z|\), не изменяя его
  чётную степень. Поэтому свойство логарифма можно применить, не исключая
  отрицательные значения \(z\).
\end{notebox}

\stepheading{Шаг 4}{Применяем правило к основной задаче}

\begin{fillbox}[Применяем правило к нашей задаче]
  \[
    \lg((x^2-4)^2)=2\lg\answerblank[32mm],
  \]
  \[
    \lg((x^2-4)^4)=4\lg\answerblank[32mm].
  \]
\end{fillbox}

\begin{warningbox}[Нельзя терять часть ОДЗ]
  Запись
  \[
    \lg((x^2-4)^4)=4\lg(x^2-4)
  \]
  требует \(x^2-4>0\) и поэтому не работает на всей ОДЗ исходного выражения.
  При этом исходные логарифмы определены на всём промежутке \(-2<x<2\);
  из общей области допустимых значений исключаются только точки \(x=\pm2\).
  Корректная запись содержит \(\lvert x^2-4\rvert\).
\end{warningbox}

\begin{thinkbox}[Карта преобразования исходного неравенства]
  Мы выяснили, как правильно преобразовать оба логарифма. Теперь вернёмся к
  исходному неравенству. Заполняйте карту постепенно по мере решения, а не
  пытайтесь завершить её сразу.
  \logchainstudentfigure
\end{thinkbox}

\section{Первая замена}

Какое выражение повторяется в преобразованном неравенстве?
\[
  \answerblank[65mm].
\]
Обозначим его через \(y\):
\[
  y=\lg|x^2-4|.
\]

\begin{fillbox}[Перепишите исходные выражения через \(y\)]
  \begin{center}
    \renewcommand{\arraystretch}{1.65}
    \begin{tabularx}{0.94\textwidth}{@{}
      >{\centering\arraybackslash}p{0.55\textwidth}
      >{\centering\arraybackslash}X@{}}
      \toprule
      \rowcolor{TableHeader}
      Исходное выражение & Через \(y\) \\
      \midrule
      \(\lg((x^2-4)^2)\) & \(\answerblank[24mm]\) \\
      \addlinespace[0.25em]
      \(\lg((x^2-4)^4)\) & \(\answerblank[24mm]\) \\
      \addlinespace[0.25em]
      \(\lg^4((x^2-4)^2)\) & \(\answerblank[30mm]\) \\
      \addlinespace[0.25em]
      \(\lg^2((x^2-4)^4)\) & \(\answerblank[30mm]\) \\
      \bottomrule
    \end{tabularx}
  \end{center}
\end{fillbox}

После подстановки:
\[
  (2y)^4-(4y)^2\ge192.
\]
Выполните вычисления:
\[
  \answerblank[16mm]y^4-\answerblank[16mm]y^2\ge192,
\]
\[
  y^4-y^2-\answerblank[14mm]\ge0.
\]

\section{Ищем вторую замену}

\begin{fillbox}[Вводим квадрат повторяющегося выражения]
  Заметим:
  \[
    y^4=\bigl(\answerblank[16mm]\bigr)^2.
  \]
  Какое выражение повторяется в неравенстве \(y^4-y^2-12\ge0\)?
  \[
    \answerblank[25mm].
  \]
  Введите замену и укажите ограничение:
  \[
    t=\answerblank[20mm],
    \qquad
    t\answerblank[10mm]0.
  \]
\end{fillbox}

При замене \(t=y^2\) получаем квадратное неравенство:
\[
  t^2-t-12\ge0.
\]
Разложите левую часть на множители:
\[
  t^2-t-12=(t-\answerblank[12mm])(t+\answerblank[12mm]).
\]

Решение квадратного неравенства:
\[
  t\in\answerblank[60mm].
\]
С учётом \(t\ge0\) остаётся:
\[
  t\ge\answerblank[14mm].
\]


\section{Возвращаемся к логарифму}

Какое множество задаёт неравенство \(y^2\ge4\)? Отметьте верный вариант:
\[
  \square\; y\ge2
  \qquad
  \square\; -2\le y\le2
  \qquad
  \square\; y\le-2\ \text{или}\ y\ge2.
\]

\begin{checkpointbox}[Не потеряйте отрицательную ветвь]
  \[
    t\ge4
    \quad\Longrightarrow\quad
    y^2\ge4
    \quad\Longleftrightarrow\quad
    |y|\ge\answerblank[10mm].
  \]
  Поэтому
  \[
    y\le\answerblank[12mm]
    \qquad\text{или}\qquad
    y\ge\answerblank[12mm].
  \]
\end{checkpointbox}

Возвращаемся к \(y=\lg|x^2-4|\):
\[
  \lg|x^2-4|\le-2
  \qquad\text{или}\qquad
  \lg|x^2-4|\ge2.
\]

\begin{thinkbox}[Почему знак неравенства сохраняется?]
  Основание десятичного логарифма равно \(10\). Так как \(10>1\), функция
  \(u\mapsto\lg u\):
  \[
    \square\; \text{возрастает}
    \qquad
    \square\; \text{убывает}.
  \]
  Поэтому при переходе от логарифма к его аргументу знак неравенства
  \choicebox\ сохраняется \qquad \choicebox\ меняется.
\end{thinkbox}

\begin{fillbox}[Переходим от логарифма к его аргументу]
  \begin{center}
    \renewcommand{\arraystretch}{1.75}
    \begin{tabularx}{0.93\textwidth}{@{}
      >{\centering\arraybackslash}p{0.40\textwidth}
      >{\centering\arraybackslash}X@{}}
      \toprule
      \rowcolor{TableHeader}
      Условие на логарифм & Условие на \(|x^2-4|\) \\
      \midrule
      \(\lg|x^2-4|\le-2\) & \(0<|x^2-4|\le\answerblank[22mm]\) \\
      \addlinespace[0.35em]
      \(\lg|x^2-4|\ge2\) & \(|x^2-4|\ge\answerblank[22mm]\) \\
      \bottomrule
    \end{tabularx}
  \end{center}
\end{fillbox}

\argregimesstudentfigure

\begin{reflectionbox}[Смысл двух ветвей]
  В первом случае \(|x^2-4|\) должно быть \answerblank[45mm],
  а во втором — \answerblank[45mm].
\end{reflectionbox}

\section{Как читать двойное неравенство с модулем}

\begin{methodintro}[Сначала упростим запись]
  Если под модулем находится сложное выражение, временно обозначьте его
  одной буквой. Положим
  \[
    u=x^2-4.
  \]
  Тогда ветвь малого аргумента принимает более простой вид:
  \[
    0<|u|\le\frac1{100}.
  \]
\end{methodintro}

\Needspace{20\baselineskip}
\begin{fillbox}[Два условия должны выполняться одновременно]
  Двойное неравенство
  \[
    0<|u|\le a,\qquad a>0,
  \]
  означает систему двух условий:
  \[
    |u|>0
    \qquad\text{и}\qquad
    |u|\le a.
  \]
  Первое условие исключает значение
  \[
    u=\answerblank[12mm].
  \]
  Второе условие равносильно двойному неравенству
  \[
    -\answerblank[12mm]\le u\le\answerblank[12mm].
  \]
  Поэтому середину промежутка нужно выколоть:
  \[
    0<|u|\le a
    \quad\Longleftrightarrow\quad
    -a\le u<0
    \quad\text{или}\quad
    0<u\le a.
  \]
\end{fillbox}

\begin{checkpointbox}[Проверьте понимание на простом примере]
  Решите неравенство
  \[
    0<|u|\le3.
  \]
  Заполните промежутки:
  \[
    u\in
    \left[\answerblank[11mm],\answerblank[11mm]\right)
    \cup
    \left(\answerblank[11mm],\answerblank[11mm]\right].
  \]
  Почему число (0) не входит в ответ, а числа \(-3\) и (3) входят?
  \writinglines{1}
\end{checkpointbox}

\begin{notebox}[Два полезных шаблона]
  При \(a>0\):
  \[
    0<|u|\le a
    \quad\Longleftrightarrow\quad
    -a\le u<0
    \quad\text{или}\quad
    0<u\le a;
  \]
  \[
    |u|\ge a
    \quad\Longleftrightarrow\quad
    u\le-a
    \quad\text{или}\quad
    u\ge a.
  \]
  Первый шаблон понадобится сейчас, второй — в ветви большого аргумента.
\end{notebox}

\section{Малый положительный аргумент}

Решаем:
\[
  0<|x^2-4|\le\frac1{100}.
\]

\begin{fillbox}[Раскрываем модуль]
  Используем \(u=x^2-4\) и первый шаблон:
  \[
    -\frac1{100}\le x^2-4<0
    \qquad\text{или}\qquad
    0<x^2-4\le\frac1{100}.
  \]
  Почему значение \(x^2-4=0\) не входит ни в один из промежутков?
  \writinglines{1}
  После прибавления \(4\):
  \[
    \frac{\answerblank[13mm]}{100}\le x^2<4
    \qquad\text{или}\qquad
    4<x^2\le\frac{\answerblank[13mm]}{100}.
  \]
\end{fillbox}

\begin{fillbox}[Переходим от \(x^2\) к \(x\)]
  \begin{center}
    \renewcommand{\arraystretch}{1.85}
    \begin{tabularx}{\textwidth}{@{}
      >{\centering\arraybackslash}p{0.30\textwidth}
      >{\centering\arraybackslash}p{0.25\textwidth}
      >{\centering\arraybackslash}X@{}}
      \toprule
      \rowcolor{TableHeader}
      Условие на \(x^2\) & Условие на \(|x|\) & Промежутки для \(x\) \\
      \midrule
      \(\dfrac{399}{100}\le x^2<4\)
        & \(\dfrac{\sqrt{399}}{10}\le|x|<2\)
        & \answerblank[57mm] \\
      \addlinespace[0.45em]
      \(4<x^2\le\dfrac{401}{100}\)
        & \(2<|x|\le\dfrac{\sqrt{401}}{10}\)
        & \answerblank[57mm] \\
      \bottomrule
    \end{tabularx}
  \end{center}
\end{fillbox}

\begin{checkpointbox}[Почему точки \(\pm2\) выколоты?]
  При \(x=\pm2\) величина \(|x^2-4|\) равна \answerblank[12mm].
  Поэтому логарифм \answerblank[40mm].
\end{checkpointbox}

\section{Большой аргумент}

Решаем:
\[
  |x^2-4|\ge100.
\]
Раскройте модуль:
\[
  x^2-4\le\answerblank[18mm]
  \qquad\text{или}\qquad
  x^2-4\ge\answerblank[18mm].
\]

\begin{thinkbox}[Одна ветвь невозможна]
  Для любого действительного \(x\):
  \[
    x^2-4\ge\answerblank[12mm].
  \]
  Поэтому условие \(x^2-4\le-100\)
  \choicebox\ возможно \qquad \choicebox\ невозможно.
\end{thinkbox}

Остаётся:
\[
  x^2\ge\answerblank[15mm]
  \quad\Longleftrightarrow\quad
  |x|\ge\answerblank[22mm].
\]
Следовательно,
\[
  x\in\answerblank[100mm].
\]

\newpage
\section{Собираем окончательный ответ}

\begin{center}
  \renewcommand{\arraystretch}{1.8}
  \begin{tabularx}{\textwidth}{@{}
    >{\centering\arraybackslash}p{0.34\textwidth}
    >{\centering\arraybackslash}X@{}}
    \toprule
    \rowcolor{TableHeader}
    Режим & Решения для \(x\) \\
    \midrule
    \(0<|x^2-4|\le10^{-2}\) & \answerblank[96mm] \\
    \addlinespace[0.55em]
    \(|x^2-4|\ge10^2\) & \answerblank[96mm] \\
    \bottomrule
  \end{tabularx}
\end{center}

\finalnumberlinestudentfigure

\begin{resultbox}[Ответ]
  \writinglines{3}
\end{resultbox}

\begin{checkpointbox}[Проверяем характерные точки]
  \begin{center}
    \renewcommand{\arraystretch}{1.55}
    \begin{tabularx}{\textwidth}{@{}
      >{\centering\arraybackslash}p{0.20\textwidth}
      >{\centering\arraybackslash}p{0.22\textwidth}
      >{\centering\arraybackslash}p{0.27\textwidth}
      >{\centering\arraybackslash}X@{}}
      \toprule
      \rowcolor{TableHeader}
      \(x\) & ОДЗ & \(|x^2-4|\) & Входит в ответ? \\
      \midrule
      \(0\) & \answerblank[14mm] & \answerblank[20mm] & \answerblank[20mm] \\
      \addlinespace[0.2em]
      \(\dfrac{\sqrt{399}}{10}\) & \answerblank[14mm] & \answerblank[20mm] & \answerblank[20mm] \\
      \addlinespace[0.2em]
      \(2\) & \answerblank[14mm] & \answerblank[20mm] & \answerblank[20mm] \\
      \addlinespace[0.2em]
      \(2\sqrt{26}\) & \answerblank[14mm] & \answerblank[20mm] & \answerblank[20mm] \\
      \bottomrule
    \end{tabularx}
  \end{center}
\end{checkpointbox}

\section{Восстанавливаем решение по краткому плану}

\begin{practicebox}[Полное решение]
  Закройте предыдущие страницы и запишите краткое связное решение. Используйте
  план только для проверки полноты:
  \begin{enumerate}
    \item расшифровка степени логарифма и ОДЗ;
    \item формулы с модулем;
    \item замены \(y=\lg|x^2-4|\) и \(t=y^2\);
    \item переход \(t\ge4\Rightarrow |y|\ge2\);
    \item обе ветви для \(|x^2-4|\);
    \item точный ответ с радикалами.
  \end{enumerate}
  \workarea{120mm}
\end{practicebox}

\begin{questionsbox}[Самопроверка]
  \begin{itemize}[label=\choicebox,leftmargin=2.4em]
    \item Я различаю \(\lg^2 A\), \(\lg(A^2)\) и \(\log_2 A\).
    \item В формулах \(\lg(z^2)=2\lg|z|\) и \(\lg(z^4)=4\lg|z|\) я сохранил модуль.
    \item Я указал \(t\ge0\).
    \item Из \(y^2\ge4\) я получил две ветви.
    \item В малой ветви я сохранил строгую положительность аргумента.
    \item Я не включил \(x=\pm2\).
  \end{itemize}
\end{questionsbox}
